\section{Problem Statement}

We consider a distributed edge caching system in which a set of edge nodes, such as WiFi access points, serve content requests generated by mobile users. Each node has limited cache capacity and may satisfy a request in one of three ways: by serving the content from its own local cache, by retrieving it from another node in the network, or by fetching it from the origin server. Since users move between nodes, request patterns observed at different nodes are not independent. Instead, mobility induces structural correlations among nodes, which makes cache placement a network-level decision problem rather than a purely local one.

We represent the system as a weighted graph
\[
G = (V, E, W),
\]
where \(V\) is the set of edge nodes, \(E\) is the set of connectivity relations, and \(W = \{w_{ij}\}\) denotes mobility-induced interaction strengths. Intuitively, a larger value of \(w_{ij}\) indicates that nodes \(i\) and \(j\) are more strongly related, for example through frequent user handoffs or overlapping mobility patterns.

The objective is to design a topology-aware cooperative caching strategy that improves content accessibility while controlling replication overhead and reducing unnecessary redundancy across strongly related nodes.

\subsection{Subproblems}

\subsubsection{Demand Modeling under Partial Observability}
The available mobility trace provides user--node association information but does not include content-level request logs. Therefore, the request process must be modeled synthetically. This creates a demand estimation problem: the workload should preserve realistic properties such as skewed popularity, spatial locality, and temporal variation, while remaining simple enough for controlled evaluation.

\subsubsection{Cooperative Cache Placement}
Cache placement decisions are coupled across nodes. If each node caches content independently, neighboring nodes may store the same items and waste storage. If all decisions are fully centralized, the coordination overhead becomes impractical. The system therefore requires a cooperative placement strategy that uses topology information without assuming unrealistic global synchronization.

\subsubsection{Access--Replication Tradeoff}
Replicating content across multiple nodes improves accessibility and lowers response time, but it also consumes storage and incurs maintenance or update cost. An effective strategy must balance the gain in access efficiency against the cost of replication.

\subsubsection{Topology-Aware Redundancy Control}
Nodes with strong interaction weights tend to experience correlated demand. Storing identical content at all such nodes may reduce effective cache diversity. A useful caching policy should therefore discourage redundant replication among strongly connected nodes while preserving accessibility.

\subsection{Problem Formulation}

We model the edge network as a weighted graph \(G = (V, E, W)\), where \(V\) is the set of edge nodes, \(E\) denotes connectivity between nodes, and \(w_{ij}\) represents the strength of mobility-induced interaction between nodes \(i\) and \(j\). Stronger edge weights indicate higher mobility correlation and therefore higher potential for cooperative caching.

Let \(\mathcal{C}\) denote the content catalog. For each node \(i \in V\) and content \(c \in \mathcal{C}\), we define a binary placement variable:
\[
x_{i,c} =
\begin{cases}
1, & \text{if content } c \text{ is cached at node } i,\\
0, & \text{otherwise.}
\end{cases}
\]

Let \(P_{i,c}\) denote the probability that node \(i\) requests content \(c\), and let \(f_c\) denote the replication-related cost of content \(c\), which may reflect update frequency, storage burden, or a combined maintenance factor.

Let \(d(i,j)\) be the shortest-path distance between nodes \(i\) and \(j\), and let \(L(i,j)\) denote the latency of retrieving content from node \(j\) to node \(i\). We assume that \(L(i,j)\) is non-decreasing with graph distance. Let \(L_{\text{local}}\) denote the latency of serving content from the local cache, and let \(L_{\text{origin}}\) denote the latency of fetching it from the origin server.

To simplify the notation, we define the best non-local retrieval latency for content \(c\) at node \(i\) under placement \(x\) as
\[
A_{i,c}(x)
=
\min\!\left(
L_{\text{origin}},
\;
\min_{\substack{j \in V \\ j \neq i \\ x_{j,c}=1}} L(i,j)
\right).
\]
If no other node caches content \(c\), then \(A_{i,c}(x)=L_{\text{origin}}\).

The expected access cost is defined as
\[
C_{\text{access}}(x)
=
\sum_{i \in V}\sum_{c \in \mathcal{C}}
P_{i,c}
\left[
x_{i,c}L_{\text{local}}
+
(1-x_{i,c})A_{i,c}(x)
\right].
\]
This term captures the average latency experienced by requests under the current cache placement.

The total replication cost is given by
\[
C_{\text{rep}}(x)
=
\sum_{i \in V}\sum_{c \in \mathcal{C}}
\gamma f_c x_{i,c},
\]
where \(\gamma > 0\) scales the importance of replication overhead. This term penalizes storing too many replicas, especially for contents that are expensive to maintain or update.

To discourage unnecessary duplication among strongly related nodes, we define the topology-aware redundancy term:
\[
C_{\text{red}}(x)
=
\sum_{c \in \mathcal{C}}
\sum_{\substack{i,j \in V \\ i<j}}
w_{ij}\,x_{i,c}x_{j,c}.
\]
This term becomes large when the same content is placed at multiple nodes that have strong mobility-induced interaction, thus reducing effective cache diversity.

The overall objective is to jointly minimize access cost, replication cost, and redundancy. We formulate the topology-aware cooperative caching problem as
\[
\min_{x}
\quad
C_{\text{access}}(x)
+
\lambda C_{\text{rep}}(x)
+
\beta C_{\text{red}}(x),
\]
where \(\lambda > 0\) and \(\beta > 0\) are tradeoff parameters. Specifically, \(\lambda\) controls the relative penalty assigned to replication overhead, while \(\beta\) controls the strength of redundancy suppression among strongly connected nodes.

This formulation captures the central design tradeoff of the system. A placement that aggressively replicates content may reduce access latency but will incur larger replication and redundancy costs. Conversely, a placement that overly avoids replication may reduce storage overhead but increase access latency due to more frequent remote retrievals. The optimal caching policy therefore balances these competing effects under storage and communication constraints.

\subsection{Constraints}

The optimization problem is subject to several practical constraints.

\paragraph{Cache Capacity Constraint.}
Each edge node has limited storage and therefore cannot cache the entire content catalog. If \(K_i\) denotes the cache capacity of node \(i\), then
\[
\sum_{c \in \mathcal{C}} x_{i,c} \le K_i,
\qquad \forall i \in V.
\]

\paragraph{Binary Placement Constraint.}
Caching is a discrete decision, so
\[
x_{i,c} \in \{0,1\},
\qquad \forall i \in V,\; \forall c \in \mathcal{C}.
\]

\paragraph{Latency Ordering Constraint.}
Serving content locally is assumed to be cheaper than retrieving it from another node, and origin access is the most expensive:
\[
L_{\text{local}} < L(i,j) \le L_{\text{origin}},
\qquad \forall i,j \in V,\; i \neq j.
\]

\paragraph{Topology-Dependent Retrieval Constraint.}
Inter-node retrieval cost depends on the graph structure through \(L(i,j)\), which is determined by path distance and connectivity conditions. This ensures that the topology directly influences the cost of cooperative content access.


\begin{figure}[t]
\centering
\begin{tikzpicture}[
    node distance=2.0cm and 2.4cm,
    ap/.style={
        draw,
        rounded corners,
        minimum width=1.8cm,
        minimum height=1.0cm,
        align=center,
        font=\small
    },
    cache/.style={
        draw,
        rectangle split,
        rectangle split parts=2,
        rounded corners,
        minimum width=1.8cm,
        minimum height=0.9cm,
        align=center,
        font=\scriptsize
    },
    origin/.style={
        draw,
        ellipse,
        minimum width=2.2cm,
        minimum height=1cm,
        align=center,
        font=\small
    },
    edge/.style={thick},
    req/.style={-{Latex[length=2.5mm]}, thick},
    coop/.style={-{Latex[length=2.5mm]}, dashed, thick},
    note/.style={font=\scriptsize, align=center}
]

% Nodes
\node[cache] (ap1) {AP1 \nodepart{second} Cache: $c_1,c_3$};
\node[cache, right=of ap1] (ap2) {AP2 \nodepart{second} Cache: $c_2,c_4$};
\node[cache, below=of ap1] (ap3) {AP3 \nodepart{second} Cache: $c_1,c_5$};
\node[cache, below=of ap2] (ap4) {AP4 \nodepart{second} Cache: $c_3,c_6$};

% Origin
\node[origin, right=3.2cm of ap2] (origin) {Origin Server};

% Mobility-induced graph edges
\draw[edge] (ap1) -- node[above, font=\scriptsize] {$w_{12}$} (ap2);
\draw[edge] (ap1) -- node[left, font=\scriptsize] {$w_{13}$} (ap3);
\draw[edge] (ap2) -- node[right, font=\scriptsize] {$w_{24}$} (ap4);
\draw[edge] (ap3) -- node[below, font=\scriptsize] {$w_{34}$} (ap4);
\draw[edge] (ap1) -- node[sloped, above, font=\scriptsize] {$w_{14}$} (ap4);

% Request arrows
\draw[req] ([xshift=-0.9cm]ap1.west) -- node[above, font=\scriptsize] {request $c_2$} (ap1.west);

% Cooperative retrieval
\draw[coop, bend left=15] (ap1.north east) to node[above, font=\scriptsize] {neighbor fetch} (ap2.north west);

% Origin fetch
\draw[req, bend left=10] (ap3.east) to node[above, font=\scriptsize] {origin fetch} (origin.west);

% Labels / notes
\node[note, below=0.3cm of ap3] {Local hit if content is in \\ node cache};
\node[note, below=0.3cm of ap4] {Otherwise retrieve from \\ another node or origin};

\end{tikzpicture}
\caption{Topology-aware cooperative caching problem. Each edge node has limited cache capacity and participates in a mobility-induced graph. A request may be served locally, cooperatively from another node, or from the origin server, depending on cache placement and graph connectivity.}
\label{fig:problem_model}
\end{figure}